\subsection{Why a relation graph rather than a moving overlay?}
The proposed graph adds a relation layer between two commonly separated data sources. TSP data describe anomalous geological conditions, but they do not specify which TBM components are geometrically related to the anomalous voxels during excavation. TBM monitoring records describe machine responses, but they do not identify the spatial source of the response in the surrounding geological field. The graph connects these perspectives by representing candidate spatial relations between rock voxels and component-labelled TBM surface nodes.

This contribution is different from producing another monitoring indicator. The component-level geological anomaly index is derived from candidate rock--TBM relations and therefore preserves component identity, relation-screening attributes, edge lifecycle and edge provenance. These properties are lost in global geological summaries or pooled TBM summaries. They are essential for shield jamming because the abnormal response is component-specific rather than a machine-wide phenomenon.

\subsection{What the event evidence supports}
The event-centred evidence supports four representation claims. First, the graph reconstructs a moving relation-support process in which anomalous rock voxels become increasingly associated with shield components as the TBM approaches the confirmed event. Second, the component-level index provides more specific evidence than global or pooled geological summaries. Third, the relation pattern can be checked against parameter sensitivity and ablation settings. Fourth, the interpretation can be traced back to specific candidate edges and their geological and geometric attributes.

The evidence should be read at this representation level. The spatial relation weight is a geometry-compatibility weight, not a measurement of ground--machine contact force, shield pressure, friction or risk probability. The component-level geological anomaly index is an edge-weighted anomaly summary over a screened relation support; it is not a calibrated failure probability. Response-alignment diagnostics can support consistency with the event, but they should not be presented as causal proof of a unique failure mechanism.

\subsection{Implications for spatial data modelling in TBM excavation}
The proposed representation treats TBM excavation as a moving spatial relation process. Rock voxels, TBM surface nodes, candidate rock--TBM relations, edge attributes and edge-level provenance form a graph-structured description of the excavation scene. This representation can support spatial queries such as which anomalous voxels are related to a given shield component, which candidate edges contribute most to a component-level index, and how the support of rock--TBM relations changes across excavation steps.

The same relation layer can be linked to future prediction modules without being absorbed into a black-box model. A monitoring-based predictor could use $I_c(t)$, $\Delta I_c(t)$, $A_c(t)$ or edge-provenance summaries as structured inputs, while the graph would retain the entity and edge records needed for spatial interpretation. This separation is important because operational prediction and spatial provenance answer different questions.

\subsection{Limitations and future work}
The main limitation is the event-centred nature of the evaluation. A confirmed shield-jamming event provides a strong engineering anchor, but one event cannot establish universal prediction ability. Future work should test the same representation on additional shield-jamming, non-jamming and mixed-anomaly intervals. The second limitation is the treatment of TSP-derived anomaly scores as deterministic voxel attributes. Future work should propagate TSP uncertainty through $q_i$, test alternative geological anomaly scores and examine how uncertainty affects top-edge provenance. The third limitation is parameter dependence in the graph-construction stage. The sensitivity design proposed in this paper should be expanded into a multi-case robustness analysis once more events become available.
